\section{Lecture 4 (January 17, 2025): Representations as comodules}\label{lec4}
\noindent\textit{Scribe: Charlie Magland}

\subsection{The representation--comodule dictionary}
Let $G$ be an affine group scheme over a field $k$, and put $A=k[G]$.
We prove Theorem~\ref{rep-comodule-dictionary}.

\begin{proof}
Given a right coaction $\rho(v)=\sum v_0\otimes v_1$, define, for
$g:A\to R$,
\[
 r_R(g)(v\otimes a)=\sum v_0\otimes g(v_1)a.
\]
This map is $R$-linear. Coassociativity and the definition of the product
in $G(R)$ give $r_R(g)r_R(h)=r_R(gh)$; counitality gives
$r_R(e)=\id$. In particular $r_R(g^{-1})$ is its inverse. Composition
of algebra maps $R\to R'$ proves naturality.

Conversely, suppose $r_R$ is such a natural action. Evaluate at the
universal point $u=\id_A\in G(A)$ and define
\[
 \rho(v)=r_A(u)(v\otimes1).
\]
For every $g:A\to R$, naturality gives
$r_R(g)(v\otimes1)=(\id_V\otimes g)\rho(v)$, so this coaction, if
valid, recovers the action. Over $A\otimes A$, the two points
$\iota_l(a)=a\otimes1$ and $\iota_r(a)=1\otimes a$ multiply to
$\Delta$. Applying the action law and naturality to this identity gives
$(\rho\otimes\id)\rho=(\id\otimes\Delta)\rho$. Evaluating at
$\epsilon:A\to k$, the identity of $G(k)$, gives counitality.
The displayed formulas show that the constructions are inverse.
Notice that $u=\id_A$ is the universal point, not generally the identity
point $u_A\epsilon$ of $G(A)$.
\end{proof}

\begin{definition}
A morphism of right $A$-comodules $f:M\to N$ is a linear map satisfying
\[
 \rho_Nf=(f\otimes\id_A)\rho_M.
\]
By the formulas above, this is equivalent to $G(R)$-equivariance after
extension of scalars to every $R$.
\end{definition}

\begin{example}
The \emph{regular comodule} is $A\xrightarrow{\Delta}A\otimes A$.
The \emph{trivial comodule} is $k\to k\otimes A$, $1\mapsto1\otimes1$;
it uses the unit $k\to A$, not the counit $A\to k$.
\end{example}

\subsection{The abelian category of comodules}
The next constructions require only that $A$ be a coalgebra over $k$.
A subspace $N\subseteq M$ is a \emph{subcomodule} if
$\rho_M(N)\subseteq N\otimes A$. The quotient then has the well-defined
coaction
\[
 \overline\rho(m+N)=\sum(m_0+N)\otimes m_1.
\]
Coassociativity and counitality follow from those on $M$.

\begin{theorem}
The category of right $A$-comodules is a $k$-linear abelian category.
Kernels, cokernels, images and direct sums are computed on the underlying
vector spaces. The same holds for the full subcategory of finite-dimensional
comodules, with finite direct sums. The forgetful functors are exact.
\end{theorem}
\begin{proof}
For a comodule map $f:M\to N$, the identity
$(f\otimes\id_A)\rho_M=\rho_Nf$ and exactness of tensoring over a field
give
\[
 \rho_M(\ker f)\subseteq\ker(f\otimes\id_A)=\ker(f)\otimes A.
\]
The image is also a subcomodule, and the ordinary quotient $N/\im f$
inherits a coaction. These constructions have the kernel and cokernel
universal properties because their factorizations are linear maps whose
coaction identities descend from that of $f$. Moreover the vector-space
isomorphism $M/\ker f\simeq\im f$ respects coactions.

The direct sum has componentwise coaction. For a finite sum $M_1\oplus M_2$,
the usual inclusions $i_j$ and projections $p_j$ are comodule maps and satisfy
\[
 p_ji_j=\id_{M_j},\qquad p_ki_j=0\ (k\ne j),\qquad
 i_1p_1+i_2p_2=\id_{M_1\oplus M_2}.
\]
The zero object and these biproducts make the category additive; comodule
maps form linear subspaces of ordinary Hom spaces. The kernel, cokernel and
image--coimage statements therefore prove that it is abelian. All these
constructions preserve finite-dimensionality when the inputs are finite-dimensional.
\end{proof}

\subsection{Local finiteness}
\begin{theorem}\label{local-finiteness-comodules}
Every element of an $A$-comodule belongs to a finite-dimensional subcomodule.
Consequently every comodule is the directed union of its finite-dimensional
subcomodules.
\end{theorem}
\begin{proof}
Write $\rho(m)=\sum_{i=1}^r m_i\otimes a_i$ with the $a_i$ linearly
independent, and let $N=\Span\{m_1,\ldots,m_r\}$. Choose linear
functionals $\lambda_j\in A^*$ with $\lambda_j(a_i)=\delta_{ij}$.
Coassociativity gives
\[
 \sum_i\rho(m_i)\otimes a_i=\sum_i m_i\otimes\Delta(a_i).
\]
Applying $\id_M\otimes\id_A\otimes\lambda_j$ shows
$\rho(m_j)\in N\otimes A$. Counitality gives
$m=\sum_i\epsilon(a_i)m_i\in N$. Thus $N$ is the required subcomodule.
Finite sums of finite-dimensional subcomodules are finite-dimensional,
which proves the directed-union assertion.
\end{proof}

\begin{definition}
An affine group scheme is \emph{algebraic} if its coordinate algebra is
finitely generated as a $k$-algebra.
\end{definition}

\begin{corollary}\label{affine-inverse-limit}
Every affine group scheme over $k$ is an inverse limit of algebraic affine
group schemes.
\end{corollary}
\begin{proof}
It suffices to show that a commutative Hopf algebra $A$ is the filtered union
of its finitely generated Hopf subalgebras. Given $a\in A$, choose a
finite-dimensional right subcomodule $V\subseteq A$ containing $a$.
In a basis $v_1,\ldots,v_n$ write
$\Delta(v_j)=\sum_i v_i\otimes c_{ij}$. The coaction identities give
\[
 \Delta(c_{ij})=\sum_l c_{il}\otimes c_{lj},\qquad
 \epsilon(c_{ij})=\delta_{ij}.
\]
Applying $\epsilon\otimes\id$ to $\Delta(v_j)$ shows that every $v_j$
is a linear combination of the $c_{ij}$. Hence the algebra generated by
the $c_{ij}$ and $S(c_{ij})$ contains $a$. It is stable under
comultiplication and counit, and also under $S$, since $S^2=\id$ in a
commutative Hopf algebra (inversion squared is the identity). It is
therefore a finitely generated Hopf subalgebra. Applying the same argument
to finite sets makes these subalgebras a directed family. Taking spectra
reverses the resulting filtered union to an inverse limit.
\end{proof}
The subalgebras here are finitely generated as algebras, not necessarily
finite-dimensional as vector spaces; $k[x]$ already illustrates the distinction.

\subsection{Tensor products}
For any bialgebra $A$, the tensor product of right comodules has coaction
\[
 \rho_{M\otimes N}(m\otimes n)
   =\sum m_0\otimes n_0\otimes m_1n_1.
\]
Equivalently, apply $\rho_M\otimes\rho_N$, interchange the middle factors
using $\tau:A\otimes N\to N\otimes A$, then multiply the two $A$ factors.
The bialgebra axioms imply the coaction identities. The unit object is the
trivial comodule $k$, and the ordinary associativity and unit maps are
comodule maps. When $A$ is commutative, the flip
$M\otimes N\to N\otimes M$ is colinear, since $m_1n_1=n_1m_1$.
Thus representations of affine group schemes have a symmetric monoidal
structure. For a general Hopf algebra the ordinary flip need not be colinear.
